\emph{Översikt}
Den här övningen låter dig träna på att skriva program med ett grafiskt
gränssnitt.  Vi börjar med en klocka: först i terminalen, med en slinga
som skriver, väntar och skriver igen, och sedan i ett fönster.
Skillnaden mellan de två
lösningarna är övningens röda tråd, och den återkommer i en timer som du
ställer in och startar, och i ett ritprogram där du väljer färg och
ritar med musen.  Varje uppgift följs av ett lösningsförslag, men
lösningen är inte poängen: försök själv först, och läs sedan förslaget
som ett andra svar att jämföra med ditt eget.

\emph{Lärandemål}
Efter den här övningen ska studenten kunna:
% Lo08 (design interactive programs)
\begin{restatable}{lo}{OvningGraphicsLOEventLoop}%
  \label{OvningGraphicsLOEventLoop}%
  Förklara hur ett grafiskt program drivs av en händelseslinga och vad en
  \foreignlanguage{english}{callback}-funktion är.
\end{restatable}
% Lo05 (design and present user friendly output),
% Lo08 (design interactive programs)
\begin{restatable}{lo}{OvningGraphicsLOWidgets}%
  \label{OvningGraphicsLOWidgets}%
  Konstruera ett enkelt grafiskt gränssnitt med
  \mintinline{python}{tkinter} med fönster, etiketter, textfält och
  knappar.
\end{restatable}
% Lo08 (design interactive programs)
\begin{restatable}{lo}{OvningGraphicsLOBind}%
  \label{OvningGraphicsLOBind}%
  Koppla händelser från användaren (knapptryck, tangenter, musrörelser)
  och från klockan till funktioner och metoder som uppdaterar
  gränssnittet.
\end{restatable}
% Musrörelserna är övningens tillägg till föreläsningens formulering;
% ritprogrammet i de två sista uppgifterna kräver dem.
% Lo03 (divide a program), Lo09 (use and design composite data types)
\begin{restatable}{lo}{OvningGraphicsLOStructure}%
  \label{OvningGraphicsLOStructure}%
  Strukturera ett grafiskt program med klasser, till exempel genom att
  ärva från \mintinline{python}{tkinter}-klasser, så att gränssnitt och
  programlogik hålls isär.
\end{restatable}
% Lo06 (create flexible applications)
\begin{restatable}{lo}{OvningGraphicsLODocs}%
  \label{OvningGraphicsLODocs}%
  Hitta och använda dokumentationen för ett programbibliotek för att
  lösa ett problem biblioteket inte har något färdigt svar på.
\end{restatable}
Utöver föreläsningens mål, breddade med musrörelser, lägger övningen
till en aspekt som varje uppgift efter den första kräver:
% Lo08 (design interactive programs),
% Lo09 (use and design composite data types)
\begin{restatable}{lo}{OvningGraphicsLOState}%
  \label{OvningGraphicsLOState}%
  Lagra det ett program behöver minnas mellan två händelser i objektets
  attribut, och avgöra vad som måste sparas där.
\end{restatable}

\emph{Förkunskaper}
Studenten bör ha tagit del av föreläsningen \emph{Grafiska
användargränssnitt}, som introducerar \mintinline{python}{tkinter},
händelseslingan och det ritprogram som två av uppgifterna bygger vidare
på.  Dessutom behövs föreläsningen \emph{Klasser och objekt}: varje
lösningsförslag här är en klass som ärver från en klass i
\mintinline{python}{tkinter}, och \mintinline{python}{super()} anropas i
varje konstruktor.  En av uppgifterna läser tal som användaren skrivit
in och måste hantera att det inte är ett tal, vilket görs med
\mintinline{python}{try} och \mintinline{python}{except} från
föreläsningen \emph{Inmatning och felhantering}.  Tupler, villkor och
upprepningar används genomgående.

\emph{Läsning}
Pythons egen dokumentation för \mintinline{python}{tkinter} beskriver
varje komponent och varje metod som används
här.\footnote{%
  \url{https://docs.python.org/3/library/tkinter.html} (hämtad
  2026-09-17).}
Samma text finns i terminalen: \texttt{pydoc3 tkinter.Tk},
\texttt{pydoc3 tkinter.Canvas} och så vidare.
